\section{Walk expansions in traced domains}\label{s:walks}

This section and the next prove Hypothesis~(P) from three hypotheses, (H1), (H2) and (H3$'$), which are verified by computer in \S\ref{s:computer} (Theorem~\ref{t:inputs}). Here we set up the walk expansion of a pair correlation and its ingredients: traced domains, the step identity, the step ratios and their responses, and the reduction of the per-step bound to a finite list of local patterns. The recursion itself is in \S\ref{s:recursion}. Everything in these two sections concerns the continuous model of \S\ref{s:contmodel}.

\subsection{Domains}\label{s:domains}

Fix a finite $G \subset \Z^2$, let $S = S(G)$, and let $\tau$ be a pin configuration on all legs of $G$. Every site adjacent to $\Z^2 \setminus G$ lies in $S$.

\begin{definition}\label{d:domain}
A \emph{domain} on $G$ assigns to every $x \in G$ a status $\mathrm{U}$, $\mathrm{D}$ or $\mathrm{X}$, and to every site of status $\mathrm U$ or $\mathrm D$ a set of \emph{kept} bonds among the lattice bonds of $x$ inside $G$, such that the sites of $S$ have status $\mathrm U$ and keep all their bonds, and a site of status $\mathrm D$ has at most two present bonds (defined next). A bond is \emph{present} if neither end has status $\mathrm X$ and both ends keep it. The site measure is $\lambda_{\mathrm U} = \dd n$ for status $\mathrm U$ and $\lambda_{\mathrm D} = \frac35 \dd n + \frac25 \delta_0$ for status $\mathrm D$, a measure on $\Sph \cup \{0\}$; sites of status $\mathrm X$ are \emph{removed} and carry no variable. For a domain $D$ and a function $f$ of the spins, let
\begin{equation}\label{e:ZD}
Z^\tau_D[f] = \int f \prod_{\text{present bonds } \{x, y\}} (1 + n_x \cdot n_y) \prod_{\text{legs } (x, z), \ x \in S} (1 + \tau_{xz} \cdot n_x) \prod_{x \text{ not removed}} \lambda_{\mathrm{st}(x)}(\dd n_x),
\end{equation}
and let $Z^0_D[f]$ be the same integral without the leg factors. We write $Z^\cdot_D = Z^\cdot_D[1]$, $\E^\cdot_D[f] = Z^\cdot_D[f] / Z^\cdot_D$ and $\Psi_D(\tau) = Z^\tau_D / Z^0_D$, where the dot stands for $\tau$ or $0$. The class $\Dcons(G)$ consists of the domains in which every site of status $\mathrm U$ keeps all its bonds, every site of status $\mathrm D$ keeps exactly two bonds, and

(C) \emph{every bond between two sites that are not removed is kept by both ends or by neither.}
\end{definition}
For tracing we also use an \emph{auxiliary spin network} on $G$. It has
the same statuses, site measures, and kept-bond data as a domain, with at
most two present bonds at each $\mathrm D$ site, but the sites of $S$
may keep arbitrary subsets of their bonds; they still have status
$\mathrm U$. It also carries a scalar mass $\mu>0$. Its partition
function is $\mu$ times the right-hand side of \eqref{e:ZD}, with all
the original leg factors retained, and its free partition function omits
those leg factors. A domain is identified with the corresponding network
of mass $1$. The definitions of $\E^\cdot$, $\Psi$ and trails extend to
these networks. Whenever membership in $\Dcons(G)$ is asserted, it refers
to the underlying data with mass reset to $1$; this does not change
normalized expectations. In particular, $\Dcons(G)$ retains its original
protected-boundary requirements. The positivity, conditional independence
and conditional-mean facts below also hold for auxiliary spin networks.

The domains $G \setminus Y$ of \S\ref{s:contmodel}, with $Y \subset G \setminus S$, are the domains of $\Dcons(G)$ with statuses $\mathrm U$ and $\mathrm X$ only. A site of status $\mathrm D$ will arise from a site through which a trail has passed once; the measure $\lambda_{\mathrm D}$ replaces the bonds used by the trail.

The elementary facts of \S\ref{s:contmodel} extend to domains. The moments of $\lambda_{\mathrm D}$ of positive degree are $\frac35$ times those of $\dd n$, so free positivity holds: $\E^0_D[f] \ge 0$ for every product $f$ of scalar products $n_x \cdot n_y$. Non-adjacent sites are conditionally independent given the other spins. If a site has at most three present bonds and no leg, then by Lemma~\ref{l:35} its conditional mean given the other spins has norm at most $3/5$; for status $\mathrm D$ the conditional law is a mixture of such a law and $\delta_0$, and a neighbouring spin equal to $0$ is a degenerate pin.

\subsection{The step identity and the pair walk identity}\label{s:stepid}

\begin{lemma}\label{l:stepid}
Let $a, u_1, \dots, u_k \in \R^3$. If $k \le 4$, then
\begin{equation*}
\int (a \cdot n) \prod_{i = 1}^k (1 + n \cdot u_i) \, \dd n = \sum_{i = 1}^k \frac{a \cdot u_i}{3} \int \prod_{i' \neq i} (1 + n \cdot u_{i'}) \, \lambda_{\mathrm D}(\dd n),
\end{equation*}
and if $k \le 2$, the same holds with $\dd n$ replaced by $\lambda_{\mathrm D}$ on the left and $\lambda_{\mathrm D}$ replaced by $\frac35 \delta_0$ on the right.
\end{lemma}

\begin{proof}
By \eqref{e:moments}, for $k \le 4$ both sides of the first identity equal $\frac13 \sum_i a \cdot u_i + \frac1{15} \sum_{i < i' < i''} \bigl[ (a \cdot u_i)(u_{i'} \cdot u_{i''}) + (a \cdot u_{i'})(u_i \cdot u_{i''}) + (a \cdot u_{i''})(u_i \cdot u_{i'}) \bigr]$; on the right we used $\int \prod_{j \le 3} (1 + n \cdot u_j) \, \lambda_{\mathrm D}(\dd n) = 1 + \frac15 \sum_{j < j'} u_j \cdot u_{j'}$. For $k \le 2$ both sides of the second identity equal $\frac15 \sum_i a \cdot u_i$.
\end{proof}

A \emph{trail} in a domain or auxiliary spin network $D$ from $u$ is a
sequence $w=(v_0,\dots,v_m)$ with $v_0=u$ such that
$\{v_{i-1},v_i\}$ are distinct present bonds of $D$; $|w|=m$ is its
length. For $T\subset G$ let $\mathcal T_D(u,T)$ be the set of trails
with $v_m\in T$ and $v_i\notin T$ for $i<m$. All tracing below starts
at a site with at most three present bonds and has no departures from
$S$. The \emph{traced spin network} $D;[w]$ is obtained step by step:
at step $j$, the bond $\{v_{j-1},v_j\}$ is dropped by both ends, and
the status of $v_{j-1}$ changes from $\mathrm U$ to $\mathrm D$
(keeping its remaining bonds), or from $\mathrm D$ to $\mathrm X$.
The mass is unchanged in the first case and is multiplied by $3/5$ in
the second case, the mass of $(3/5)\delta_0$. All original leg factors
are unchanged. Thus the final mass is
$\mu(D)(3/5)^{N_{\mathrm D}(w)}$, where $N_{\mathrm D}(w)$ counts
the departures from status $\mathrm D$ during this trace. We write
$w_j=(v_0,\dots,v_j)$. A stopped trace can have a terminal site in
$S$ with its arrival bond dropped; it is an auxiliary spin network,
and need not be a domain of Definition~\ref{d:domain}.

\begin{proposition}\label{p:pairwalk}
Let $D$ be a domain or auxiliary spin network on $G$, let
$u\in G\setminus S$ be a site that is not removed and has at most
three present bonds, let $v\ne u$ be a site that is not removed, and
$T=S\cup\{v\}$. Then, with the dot standing for $\tau$ or $0$,
\begin{equation}\label{e:pairwalk}
\E^\cdot_D[n_u \cdot n_v] = \sum_{w \in \mathcal T_D(u, T)} 3^{-|w|} \, R^\cdot_D(w) \, \E^\cdot_{D;[w]}[n_{v_{|w|}} \cdot n_v], \qquad R^\cdot_D(w) = \frac{Z^\cdot_{D;[w]}}{Z^\cdot_D},
\end{equation}
where for $v_{|w|} = v$ the last factor is $\E^\cdot_{D;[w]} [|n_v|^2]$, which equals $1$ if $v$ has status $\mathrm U$ and $3/5$ if it has status $\mathrm D$ with one present bond in $D$.
\end{proposition}

\begin{proof}
Write $n_u \cdot n_v = \sum_{\beta = 1}^3 n_u^\beta n_v^\beta$ and fix $\beta$. In $Z^\cdot_D[n_u^\beta n_v^\beta]$, integrate the spin of $u$ using Lemma~\ref{l:stepid} with $a = e_\beta$; the site $u$ carries no leg because $u \notin S$, and the factor $n_v^\beta$ does not involve $n_u$. This gives a sum over the present bonds $\{u, v_1\}$ of $\frac13 Z^\cdot_{D;[(u, v_1)]}[n_{v_1}^\beta n_v^\beta]$. If $v_1 \notin T$ we repeat the step at $v_1$: it carries no leg, a site of status $\mathrm U$ has at most three remaining factors after the arrival bond is dropped, and a site of status $\mathrm D$ at most one. The procedure stops at the first visit of $T$, and a step from a site without remaining bonds contributes $0$. Every trail is finite, so after finitely many steps we obtain $Z^\cdot_D[n_u^\beta n_v^\beta] = \sum_w 3^{-|w|} Z^\cdot_{D;[w]}[n_{v_{|w|}}^\beta n_v^\beta]$. Summing over $\beta$ and dividing by $Z^\cdot_D$ gives \eqref{e:pairwalk}. If $v_{|w|}=v$, the last bracket is $|n_v|^2$. If $v$ has
status $\mathrm D$ with one present bond in $D$, that bond was the
arrival bond, so after arrival $n_v$ is independent of the rest with
law $\lambda_{\mathrm D}$, and $\E|n_v|^2=3/5$.
\end{proof}

\subsection{Step ratios and their responses}\label{s:steps}

For $w\in\mathcal T_D(u,T)$ as in Proposition~\ref{p:pairwalk}, we have
$R^\cdot_D(w)=\prod_{j=1}^{|w|}\rho^\cdot_j$ with
$\rho^\cdot_j=Z^\cdot_{D;[w_j]}/Z^\cdot_{D;[w_{j-1}]}$.
Let $y=v_{j-1}$ be the site left at step $j$, $o=v_j$ the
out-neighbor, and $E'$ the auxiliary spin network obtained from
$D;[w_{j-1}]$ by removing $y$, retaining its scalar mass. Step $j$ is
a \emph{revisit} if $y$ has status $\mathrm D$ in $D;[w_{j-1}]$ and
a \emph{first visit} otherwise. Integrating out $y$, we find:
\begin{itemize}
\item for a revisit with at most one present bond at departure,
$\rho^\tau_j=\rho^0_j=3/5$, since
$\int(1+n\cdot n_o)\lambda_{\mathrm D}(\dd n)=1$.
The only possible two-bond $\mathrm D$ departure is the initial step.
If its present neighbors are $f,o$, then
\[
\rho^\cdot_1=
\frac{3/5}{1+\frac15\E^\cdot_{E'}[P_{fo}]}\le\frac34.
\]
Indeed, every later departure has already lost its arrival bond;
a $\mathrm D$ site has at most two present bonds before that arrival;
\item for a first visit with remaining neighbours $f, g, o$, $1/\rho^\cdot_j = \E^\cdot_{E'}[I] / \E^\cdot_{E'}[J]$ with
\begin{equation}\label{e:IJ}
I = 1 + \tfrac13 (P_{fg} + P_{fo} + P_{go}), \qquad J = 1 + \tfrac15 P_{fg}, \qquad P_{ab} = n_a \cdot n_b ;
\end{equation}
\item for a first visit with remaining neighbours $f, o$, the same with $I = 1 + \frac13 P_{fo}$ and $J = 1$; with the single remaining neighbour $o$, $\rho^\cdot_j = 1$.
\end{itemize}
The first visit is \emph{straight} if $o - y = y - v_{j - 2}$ and a \emph{turn} otherwise; for $j = 1$ the role of $v_{-1}$ is played by a designated neighbour $c$ of $u$, which in \S\ref{s:recursion} is the removed centre of the pair.

\begin{lemma}\label{l:response}
Let step $j$ be a first visit at $y \notin S$, and let $r_{ab} = \E^\tau_{E'}[P_{ab}] - \E^0_{E'}[P_{ab}]$. With three remaining neighbours,
\begin{equation}\label{e:lambda}
\lambda_j := \log \rho^0_j - \log \rho^\tau_j = \alpha_j r_{fg} + \beta_j (r_{fo} + r_{go}), \qquad 0 < \alpha_j \le \bar\alpha, \quad 0 < \beta_j \le \bar\beta,
\end{equation}
where $\bar \alpha = 0.18094$ and $\bar \beta = 0.36751$. With two remaining neighbours, $\lambda_j = \beta'_j r_{fo}$ with $0 < \beta'_j \le 0.35510 \le \bar\beta$; with one, $\lambda_j = 0$. Moreover $\rho^0_j \le 1$ and $\rho^\tau_j \le 12/5$.
\end{lemma}

\begin{proof}
For $x, x' > 0$ let $\mathrm L(x, x') = (x - x')/(\log x - \log x')$ be the logarithmic mean ($\mathrm L(x, x) = x$); it is increasing in both arguments. By the mean value form of $\log x - \log x'$,
\begin{equation*}
\lambda_j = \frac{\E^\tau_{E'}[I] - \E^0_{E'}[I]}{\mathrm L_I} - \frac{\E^\tau_{E'}[J] - \E^0_{E'}[J]}{\mathrm L_J}, \qquad \mathrm L_I = \mathrm L(\E^\tau_{E'}[I], \E^0_{E'}[I]), \quad \mathrm L_J = \mathrm L(\E^\tau_{E'}[J], \E^0_{E'}[J]),
\end{equation*}
which is \eqref{e:lambda} with $\alpha_j = 1/(3 \mathrm L_I) - 1/(5 \mathrm L_J)$ and $\beta_j = 1/(3 \mathrm L_I)$. Since $y \notin S$, its neighbours are at distance at least $2$ from $\Z^2 \setminus G$ and carry no legs; in $E'$ each has at most three present bonds. They are pairwise non-adjacent, hence conditionally independent given the other spins, with conditional means $m_f, m_g, m_o$ of norm at most $3/5$. Since $\sum_{\text{pairs}} m_a \cdot m_b = \frac12 (|\sum m_a|^2 - \sum |m_a|^2) \in [-\frac{27}{50}, \frac{27}{25}]$, conditioning on the other spins gives $\E^\cdot_{E'}[I] \in [0.82, 1.36]$ and $\E^\cdot_{E'}[J] \in [0.928, 1.072]$. By free positivity, $\E^0_{E'}[I] \ge 1$ and $\E^0_{E'}[J] \ge 1$. By monotonicity of $\mathrm L$, $\mathrm L_I \ge \mathrm L(0.82, 1) = 0.907025\ldots$ and $\mathrm L_J \in [\mathrm L(0.928, 1), 1.072]$ with $\mathrm L(0.928, 1) = 0.963551\ldots$ Hence
\begin{equation*}
\beta_j \le \frac{1}{3 \mathrm L(0.82, 1)} = 0.3675017\ldots, \qquad \alpha_j \le \frac{1}{3 \mathrm L(0.82, 1)} - \frac{1}{5 \cdot 1.072} = 0.1809346\ldots, \qquad \alpha_j \ge \frac{1}{3 \cdot 1.36} - \frac{1}{5 \mathrm L(0.928, 1)} > 0.0375 .
\end{equation*}
With two remaining neighbours, $J = 1$ and $\E^\cdot_{E'}[I] \in [0.88, 1.12]$, so $\lambda_j = \beta'_j r_{fo}$ with $\beta'_j = 1/(3 \mathrm L_I) \le 1/(3 \mathrm L(0.88, 1)) = 0.35509\ldots$ Free positivity gives $\E^0_{E'}[I] - \E^0_{E'}[J] = \frac2{15} \E^0_{E'}[P_{fg}] + \frac13 \E^0_{E'}[P_{fo} + P_{go}] \ge 0$, that is, $\rho^0_j \le 1$, and similarly with two neighbours. Finally, pointwise $I \ge 1/2$ and $J \le 6/5$ (with two neighbours $I \ge 2/3$ and $J = 1$), so $\rho^\tau_j \le 12/5$.
\end{proof}

\begin{lemma}\label{l:exposure}
Let $w\in\mathcal T_D(u,T)$ under the hypotheses of
Proposition~\ref{p:pairwalk}. If $u$ initially has status $\mathrm D$,
assume in addition that it has at most one present bond in $D$.
Let $\Phi$ be the set of first-visit steps of $w$, and suppose that
$\rho^\tau_j\le x_j$ for $j\in\Phi$, with numbers $x_j\ge1$.
Then for every ordering $\sigma$ of $\Phi$,
\begin{equation}\label{e:exposure}
\bigl| R^\tau_D(w) - R^0_D(w) \bigr| \le \Bigl( \frac35 \Bigr)^{|w| - |\Phi|} \sum_{k = 1}^{|\Phi|} |\lambda_{\sigma(k)}| \prod_{l \le k} x_{\sigma(l)} .
\end{equation}
\end{lemma}

\begin{proof}
The revisit factors are $3/5$ in both products. For the first visits, with $a_j = \rho^\tau_j$ and $b_j = \rho^0_j$,
\begin{equation*}
\prod_{j \in \Phi} a_j - \prod_{j \in \Phi} b_j = \sum_{k} \Bigl( \prod_{l < k} a_{\sigma(l)} \Bigr) (a_{\sigma(k)} - b_{\sigma(k)}) \Bigl( \prod_{l > k} b_{\sigma(l)} \Bigr).
\end{equation*}
We have $a_j \le x_j$, $b_j \le 1$ by Lemma~\ref{l:response}, and $|a_j - b_j| \le |\log a_j - \log b_j| \max(a_j, b_j) \le |\lambda_j| x_j$.
\end{proof}

\subsection{Local patterns and Hypothesis (H1)}\label{s:patterns}

To bound $\rho^\tau_j$ for a first-visit step at $y \notin S$, we localize to the \emph{template} $\mathsf T(y) = \{ z \in \Z^2 : d(z, y) \le 2 \}$, which lies in $G$. After a lattice symmetry fixing $y = 0$ we may assume that the arrival site is $\mathsf s = (0, -1)$; the other neighbours are $\mathsf e = (1, 0)$, $\mathsf N = (0, 1)$, $\mathsf W = (-1, 0)$, the corners are $(\pm 1, \pm 1)$, and the tips are $(\pm 2, 0)$, $(0, \pm 2)$. The out-neighbour is $\mathsf N$ for a straight step and $\mathsf e$ or $\mathsf W$ for a turn. Each corner has two lattice neighbours outside $\mathsf T(y)$, and each tip has three; we call these legs of the template \emph{slots}, $20$ in total.

\begin{definition}\label{d:pattern}
A \emph{pattern} $\pi$ consists of an out-neighbor
$o\in\{\mathsf e,\mathsf N,\mathsf W\}$ and statuses and kept bonds
of the sites $\mathsf s,\mathsf e,\mathsf N,\mathsf W$ and of the
corners, with the site measures of Definition~\ref{d:domain}. The
passage site $y$ has status $\mathrm U$ and keeps precisely its three
bonds to $\mathsf e,\mathsf N,\mathsf W$, so $\{y,\mathsf s\}$ is
dropped. Every other site of status $\mathrm U$ keeps all four of
its lattice bonds, and $o$ is not removed. For unit vectors
$m=(m_\ell)$ on the slots, let $Z_{\mathrm{cur}}(\pi;m)$ be the
integral over the spins of $\mathsf T(y)$ minus the removed sites,
with the site measures of the statuses, of the product of the factors
of the present bonds inside $\mathsf T(y)$ and of
$1+m_\ell\cdot n_z$ for every slot $\ell$ at a non-removed site $z$;
the tips have status $\mathrm U$. Let
$Z_{\mathrm{next}}(\pi;m)$ be the same integral after the step,
that is, with the measure of $y$ replaced by $\lambda_{\mathrm D}$
and the bond $\{y,o\}$ dropped. Define
\begin{equation}\label{e:cpi}
c(\pi) = \inf_{m \in (\Sph)^{20}} \frac{Z_{\mathrm{cur}}(\pi; m)}{Z_{\mathrm{next}}(\pi; m)} .
\end{equation}
The class $\Pext$ consists of the patterns in which: $\mathsf s$ is removed or has status $\mathrm D$ keeping two of its three bonds other than $\{y, \mathsf s\}$; each of $\mathsf e, \mathsf N, \mathsf W$ has status $\mathrm U$, $\mathrm X$, or $\mathrm D$ keeping the bond to $y$ and one other bond; each corner has status $\mathrm U$, $\mathrm X$, or $\mathrm D$ keeping two of its four bonds; and condition (C) holds for the bonds between non-removed sites of $\mathsf T(y)$ other than the tips. (The statuses of the tips are not part of a pattern; a bond between a tip and its neighbour $z$ is present in $Z_{\mathrm{cur}}$ and $Z_{\mathrm{next}}$ if and only if $z$ is not removed and keeps it.) A pattern is straight or a turn according to $o$.
These are local covered spin integrals, not domains with the
protected-boundary condition on $S(G)$. The tips are the covering
$\mathrm U$ tips of Lemma~\ref{l:local}; condition (C) is not imposed
between their original global kept sets and the other template sites.
Both external slot factors at each non-removed corner are included in
the covered integrals, even when an actual external bond is absent.
Such missing factors are filled by the convex-averaging enlargement
in the proof of Lemma~\ref{l:local}.
\end{definition}

Up to the reflection $x_1 \mapsto -x_1$ and after removing duplicates that give the same function in \eqref{e:cpi} (a corner leg that is dropped becomes a slot, see the proof below, and a site without present bonds contributes a factor that does not depend on the other spins of the template and cancels in the ratio), the class $\Pext$ has $7{, }962$ elements, $2{, }696$ straight and $5{, }266$ turns; this count comes from a computer enumeration of the rules of Definition~\ref{d:pattern}.

\begin{hypothesis}[H1]\label{h:H1}
For constants $\cs, \ct \in (0, 1]$: $c(\pi) \ge \cs$ for every straight $\pi \in \Pext$ and $c(\pi) \ge \ct$ for every turn $\pi \in \Pext$.
\end{hypothesis}

By Theorem~\ref{t:H1} in \S\ref{s:H1}, Hypothesis~(H1) holds with $(\cs, \ct) = (0.961, 0.935)$.

\begin{lemma}\label{l:local}
Let step $j$ of a trail be a first visit at $y \notin S$, and suppose that, after the symmetry above, the statuses and kept bonds of $D;[w_{j - 1}]$ on $\mathsf T(y) \setminus \{\text{tips}\}$ form a pattern $\pi$. Then $1/\rho^\tau_j \ge c(\pi)$.
\end{lemma}

\begin{proof}
Let $\omega$ denote the spins outside $\mathsf T(y)$. Both $Z^\tau_{D;[w_{j - 1}]}$ and $Z^\tau_{D;[w_j]}$ are integrals over $\omega$ of the same nonnegative weight of the factors outside $\mathsf T(y)$, times the integrals $Z_{\mathrm{cur}}(\omega)$ and $Z_{\mathrm{next}}(\omega)$ over the template spins of the remaining factors. A ratio of two such integrals is at least $\inf_\omega Z_{\mathrm{cur}}(\omega)/Z_{\mathrm{next}}(\omega)$. We now enlarge the set over which the infimum is taken, in three steps; each can only lower it.

First, the dependence on $\omega$ is through the factors of the bonds from template sites to sites outside, each affine in an outside spin of norm at most $1$. Replacing these spins by independent vectors $m_\ell$ in the closed unit ball gives a ratio that is linear-fractional in each $m_\ell$ with positive denominator, so its infimum over the ball is attained on the sphere. A bond leaving the template that is absent contributes $1 = \frac12 (1 + n \cdot m) + \frac12 (1 - n \cdot m)$, and a ratio of averages is at least the smaller ratio; so we may assume a slot on every leg of the template. Second, consider a tip $t$ with its neighbour $z \in \mathsf T(y)$. Integrating the spin of $t$, in any status and with the actual outside spins, produces a factor $A (1 + b \cdot n_z)$ with $A \ge 0$ and $|b| \le 3/5$ by Lemma~\ref{l:35}, or a constant if the bond $\{t, z\}$ is absent. For a tip of status $\mathrm U$ with three slots the same integral gives $A(1 + b \cdot n_z)$ with $b = $ the mean of the three-factor law, and by Lemma~\ref{l:35}, rotation invariance and continuity, every $b$ with $|b| \le 3/5$ occurs. The ratio is linear-fractional and homogeneous in $(A, A b)$, so the infimum over all statuses is at least the infimum over $\mathrm U$-tips with three slots. Third, after these replacements the integrals are $Z_{\mathrm{cur}}(\pi; m)$ and $Z_{\mathrm{next}}(\pi; m)$.
\end{proof}

\begin{remark}
Lemma~\ref{l:local} applies to every first-visit step whose pattern is
in $\Pext$. Lemma~\ref{l:closure} in \S\ref{s:recursion} supplies
these patterns throughout the first expansion. If that expansion
stops at $e\in S$, its terminal trace is an auxiliary spin network
with the single kept-set defect at $e$. In the second expansion of
$T_2$, at most $12$ first-visit departures, at sites $y'$ with
$0<d(y',e)\le2$, may fail to have a covered pattern. At those
departures we use $\rho^\tau_j\le12/5$ from
Lemma~\ref{l:response}. These exceptional departures need not be
the last departures of the second trail.
\end{remark}

\subsection{Weighted trail sums and Hypothesis (H2)}\label{s:H2def}

Let $x_{\mathrm s}=1/\cs$ and $x_{\mathrm t}=1/\ct$.
For a trail $w=(v_0,\dots,v_m)$ of $\Z^2$ with a designated unit arrival direction $a$ at $v_0$, we require $v_1-v_0\ne-a$ when $m\ge1$.
The arrival direction is used to classify the first passage; its virtual incoming bond is not otherwise included among the bonds of $w$.
Let $\mathsf W(w)=\prod_{j=1}^m\omega_j$, with the empty product equal to $1$,
where $\omega_j=1/5$ if step $j$ leaves a site among $v_0,\dots,v_{j-2}$,
and otherwise $\omega_j=x_{\mathrm s}/3$ if $v_j-v_{j-1}$ equals the arrival direction ($a$ for $j=1$, $v_{j-1}-v_{j-2}$ for $j>1$), and $\omega_j=x_{\mathrm t}/3$ otherwise.

\begin{hypothesis}[H2]\label{h:H2}
There exist $C_V < \infty$ and $q < 1$ such that for all $n \ge 0$, all starting points and all arrival directions, $\sum_{|w| = n} \mathsf W(w) \le C_V q^n$, where the sum runs over the trails of $\Z^2$ of length $n$ with that start and arrival direction and with initial reversal forbidden as above.
\end{hypothesis}

\begin{lemma}\label{l:onepoint}
Assume (H1) and (H2). Let $D$ be a domain or auxiliary spin network
on $G$, let $x\in G\setminus S$ be a site that is not removed and
has at most three present bonds, with a designated arrival direction $a$ such that the bond $\{x-a,x\}$ is absent in $D$,
and suppose that all first-visit steps of every trail in
$\mathcal T_D(x,S)$, except at most $12$ steps of each trail, have
their patterns in $\Pext$. Then
\begin{equation*}
\sum_{w\in\mathcal T_D(x,S)}3^{-|w|}
\max\{R^\tau_D(w),R^0_D(w)\}
\le \Bigl(\frac{12}{5}\Bigr)^{12}
\frac{C_V}{1-q}\,q^{d(x,S)}.
\end{equation*}
\end{lemma}

\begin{proof}
Every trail of $D$ is a trail of $\Z^2$ with initial reversal forbidden for the designated arrival direction, since the bond $\{x-a,x\}$ is absent in $D$. We prove the same pointwise
bound for both dots $\tau$ and $0$. If the initial site has status
$\mathrm D$ and two present bonds, the formula above gives
$\rho^\cdot_1\le3/4$, hence
$3^{-1}\rho^\cdot_1\le1/4\le\omega_1$, since the first lattice
departure has $\omega_1\ge1/3$. Every other departure from status
$\mathrm D$ has at most one present bond, so
$3^{-1}\rho^\cdot_j=1/5$. This equals $\omega_j$ if the site was
left before by $w$, and is at most $\omega_j$ otherwise. For a
first-visit step whose pattern is in $\Pext$,
Lemma~\ref{l:local} and (H1) give
$\rho^\tau_j\le x_{\mathrm{type}}$, while
Lemma~\ref{l:response} gives $\rho^0_j\le1\le x_{\mathrm{type}}$;
thus $3^{-1}\rho^\cdot_j\le\omega_j$. At any other first-visit
step, Lemma~\ref{l:response} gives
$3^{-1}\rho^\cdot_j\le4/5\le(12/5)\omega_j$.
There are at most $12$ such steps, so pointwise
\[
3^{-|w|}\max\{R^\tau_D(w),R^0_D(w)\}
\le (12/5)^{12}\mathsf W(w).
\]
Since $|w|\ge d(x,S)$, summing (H2) over the possible lengths gives
the claim.
\end{proof}
